\section{Analysis \& Action Plan}
\label{sec:analysis}

Read against the definition of success in Section~1, the routing engine meets the goal the proposal set for it. Requests return
loops close to the requested distance (US1), with real alternatives (US2), and
with a quarter to a third of the interruptions a distance-only planner would pick
(US3), at a distance cost well inside the limit. The surface preference works
reliably (US5), GPX output is valid (US7), the system is fast once warm
(US10), and the same pipeline now serves six more metro areas (US11). The
first validation run found the result weaker at the point where the runner
meets the engine. The proposal promised that a runner could ``compare
measurable tradeoffs'' and ``make the final choice'' themselves (US4), and on
16 November the interface did not show most of the tradeoffs the engine was
making. We fixed that display on 18 November and the rerun passed. The same
day we changed how the elevation preference steers the search, which brought
US6 from 65.6\% to 90.6\% live. With the outside-application import
recorded on 21 November, all seven must-haves are met. One should-have, US8,
waits on two of the three members' browser checklists, and the one run we
have found a real gap: the hilly setting can leave the top route unchanged
without telling the runner.

We ranked issues by how much they affect what a runner sees and by whether
they break a must-have criterion. The two display gaps ranked first, so we
fixed them first. They were the only failures a runner would notice on every
request, and they undid the project's stated difference from existing tools,
which is that the tradeoffs are visible. The weak elevation preference ranked
next. It is a must-have, and a runner who asks for hills and gets the same
route as ``flat'' will conclude the setting does nothing; it was fixed second.
Deployment issues followed, because several fixes were written but not live;
the backend was redeployed on 18 November. Map-data and scaling issues rank lower. They are
real, but they mostly affect the metros beyond our primary area.

Tables~\ref{tab:actions} and~\ref{tab:actions-low} record a decision for each important finding. Owners
follow the areas of responsibility in our proposal: the interface for
Adithya, data, performance and deployment for Krish, and routing and scoring
for Rohan. Each item names how we will check that it worked. In every case
the check reruns the same script that found the problem, so a fix is not
counted until the measurement that failed passes. Priority follows the
ranking above: \emph{High} breaks a must-have or affects every runner,
\emph{Medium} affects some requests or the team's ability to operate the
system, and \emph{Low} is a measurement or documentation gap.

\begin{table}[htbp]
\centering\footnotesize
\renewcommand{\arraystretch}{1.0}% this table only
\caption{Decisions on the main findings: high and medium priority.}
\label{tab:actions}
\begin{tabularx}{\linewidth}{@{}>{\raggedright\arraybackslash}p{1.1cm} >{\raggedright\arraybackslash}p{2.4cm} >{\raggedright\arraybackslash}p{1.6cm} Y l >{\raggedright\arraybackslash}p{2.9cm}@{}}
\toprule
Priority & Finding & Decision & Action & Owner & How we verify \\
\midrule
High & US4: crossings, repetition, distance error and score not shown & Completed fix (\#29) & Restored crossings; added distance off target, repetition and score to each card and an expandable score breakdown & Rohan & Deployed rerun 18 Nov: 3,969 values, 0 mismatches, 8/8 components shown \\
High & US9: drive cost and score not shown & Completed fix (\#29) & Show driving distance and score on alternative-start cards and in the panel & Rohan & Ten-case rerun 18 Nov: both shown, 19/19 in radius \\
High & US6: hilly often ties flat & Completed fix, deployed (\#38) & Flat requests pay for climbing in the path search; a hilly credit was tried and dropped & Rohan & \texttt{acceptance.py}: US6 90.6\% live, 87.5\% served; US1, US3, US5 unchanged \\
High & US8: two members' checklist runs and all phone runs not recorded & Planned, before submission & Adithya and Rohan run Appendix~\ref{app:checklists} on desktop and phone; Krish repeats it on a phone & All three & Completed columns in Appendix~\ref{app:checklists}; any F becomes a new row here \\
High & Hilly can leave the top route unchanged with no explanation (checklist step 6) & Planned fix & Show a notice when the hilly top route equals the balanced one, as for doubled-back routes (\#32) & Rohan & Step 6 passes on rerun; US6 acceptance unchanged \\
Medium & Production backend behind \texttt{main}; manual deploys & Completed fix, deployed & Release script (\#18, \#35, \#39) that refuses a candidate reporting another commit; backend on \texttt{62e37f5} since 18 Nov & Rohan & \texttt{/api/version} equals the release commit; release checks pass on both public URLs \\
Medium & Precomputed routes miss 8.7\% of requests & Completed fix, deployed & Retry from the connected network and a wider ring (\#15) & Rohan & 600-request check $\geq$97\% after deploy \\
Medium & Live answers can be all slivers or repeat 23\% & Completed fix, deployed & Notice when shown routes double back (\#32); exclusion dropped, since no answer had a rounder alternative & Rohan & 48 park starts: identical routes, notice on the 12 affected answers \\
Medium & First request in a new city 4--23\,s & Partly fixed, deployed (\#24, \#25); rest in backlog & Pause garbage collection while loading; pre-load the tile when a start is set. Compact tile format (\#22) merged, not yet deployed & Krish & Pre-loaded first requests 0.4--1.8\,s (14 Nov); rerun after compact tiles ship \\
Medium & One precompute task crashed on a degenerate retry, blocking batch \texttt{b2} & Completed fix (\#26, \#41) & Cap the waypoint rescale; allow a merge across a checked commit & Krish & Regression test; benchmark and shape audit unchanged; \texttt{b2} merged (850,512 pools) \\
Medium & Start locations written to logs & Completed fix, deployed & Log the area name instead of coordinates or cells (\#31) & Rohan & Tests fail if a log line contains a start's coordinates or cell \\
Medium & Tests did not catch removed UI fields; no CI on PRs & Completed fix & Tests run on every pull request (\#21); browser tests now assert the US4 and US9 fields are present (\#29) & Rohan & CI blocks a PR that removes a required field \\
Medium & Accessibility: start pin has no accessible name; manual items unchecked & Planned fix & Give the map marker a name (or make it non-interactive); keyboard-only and screen-reader pass of the checklist & Adithya & Lighthouse: 0 failed audits; manual pass recorded \\
Medium & GPX files carry no elevation, so apps show 0 climb & Backlog & Write each point's elevation into the GPX & Krish & gpx.studio shows climb matching RunLine's within the terrain model's precision \\
\bottomrule
\end{tabularx}
\end{table}

\begin{table}[htbp]
\centering\footnotesize
\renewcommand{\arraystretch}{1.0}% this table only
\caption{Decisions on the main findings: low priority (continues Table~\ref{tab:actions}).}
\label{tab:actions-low}
\begin{tabularx}{\linewidth}{@{}>{\raggedright\arraybackslash}p{1.1cm} >{\raggedright\arraybackslash}p{2.4cm} >{\raggedright\arraybackslash}p{1.6cm} Y l >{\raggedright\arraybackslash}p{2.9cm}@{}}
\toprule
Priority & Finding & Decision & Action & Owner & How we verify \\
\midrule
Low & Start-to-finish card shows a score of about 5,000 & Planned fix & Show a comparable score, or none, for start-to-finish routes & Rohan & Step 10 card shows a score on the loop scale \\
Low & Place list served in one encoding to every client & Completed fix (\#34) & Turn off caching on the public site's \texttt{/api} forwarding route & Rohan & Each encoding correct on a preview; both rehearsal scripts pass against the public site \\
Low & Metros use graphs without crosswalk signals & Deferred & Rebuild with the signal-preserving builder in the cluster batches (\texttt{b2} and \texttt{b4} already use it) & Krish & Signal-node counts and benchmark per metro \\
Low & Drive distance can be below straight line & Completed fix, deployed (\#33) & Add the legs from each start onto the road network & Rohan & 37,061 park--start pairs: none below the straight line \\
Low & US1 5\% rule vs.\ the $\pm$0.25 mile band & Explicit decision & Keep the band; the 3 mile misses are deliberate trades for fewer interruptions & Rohan & Report both measures \\
Low & US2 overlap definition differs from the code & Backlog & Align the duplicate test with the shorter-route measure & Rohan & US2 served rerun \\
\bottomrule
\end{tabularx}
\end{table}

Several limits qualify every conclusion above.
\begin{itemize}
\item \textbf{One city.} The benchmark is 32 requests from eight starts in a
  single city. It was fixed before the system was tuned, but it is small, and
  the metro checks used six warm requests per area, where the 95th percentile
  is just the slowest request.
\item \textbf{Mapped interruptions only.} Signals and crossings are counted
  from OpenStreetMap \citep{osmwiki2026signals}. The crossing audit shows the
  planner counts what the map contains, not that the map is complete. Fewer
  signals is also not the same as safer, since a signal may protect a
  crossing.
\item \textbf{Self-built ground truth.} The crossing labels came from our own
  geometric labeller. We checked it against hand-labelled cases and inspected
  every disagreement, but it is not an independent survey.
\item \textbf{Elevation resolution.} Elevation comes from a 90\,m terrain
  model \citep{copernicus2022dem}, which captures hills but not short, steep
  grades.
\item \textbf{Test conditions.} Production timings come from one network
  path at one time of day, with no concurrent users.
\item \textbf{No outside users.} No runner outside the team has used
  RunLine under observation. Whether the tradeoffs, once displayed, actually
  help someone choose is untested, and so is whether our ranking matches
  what runners prefer.
\end{itemize}
